\begin{result}{F.3.14}{A precise hypersurface normalization construction}{Proved; historical identification retained}
Every smooth integral projective $d$-fold $X$ over a characteristic-zero field $k$ admits a finite birational morphism
\[
 \nu:X\longrightarrow Y\subset\mathbf P_k^{d+1}
\]
onto an integral hypersurface, and $\nu$ identifies $X$ with the normalization of $Y$.
This is an actual geometric construction relevant to the proposed use of singular hypersurfaces. The pullback on ordinary cohomology need not be an isomorphism: the explicit nodal cubic below has $H^1(Y(\C),\Q)=\Q$ and $H^1(X(\C),\Q)=0$, with the kernel specified by its normalization sequence.
\end{result}
\begin{proof}
The birational projection statement is explicitly present in the original EGA~V prenotes, Corollary~14.6, with its preceding Proposition~14.5 and the following warning that the projection need not be an immersion \cite[original scan p.89]{AuditGrothendieckEGAVOriginal}. We give the construction and normalization proof in the stated characteristic-zero scope.

Choose a very ample bundle $L$ and a projective embedding defined by a basis of $H^0(X,L)$. Because $k$ is infinite, choose $d+1$ sections $s_0,\ldots,s_d$ having no common zero. One can construct them successively: at each stage choose a section avoiding all the finitely many generic points of the current intersection; after $d$ cuts the intersection is finite, and the last section avoids those points. This gives
\[
 g=[s_0:\cdots:s_d]:X\longrightarrow\mathbf P_k^d,
 \qquad g^*\mathcal O(1)=L.
\]
A positive-dimensional geometric fibre would contain a projective integral curve, on which $L$ would both have positive degree and be trivial. Thus $g$ is quasi-finite; it is proper and hence finite \cite[Tag 02LS, Lemma 37.44.1]{AuditStacksFiniteNormalization}. Its image is closed of dimension $d$, so it is all of $\mathbf P^d$. Consequently
\[
 F=k(s_1/s_0,\ldots,s_d/s_0)\subset E=k(X)
\]
is a finite separable extension.

Complete the sections to a basis. Their ratios $r_j=s_j/s_0$ generate $E$ over $F$. There is a $k$-linear combination $r=\sum_j a_jr_j$ which is a primitive element. Indeed, for every pair of distinct $F$-embeddings $\tau,\tau'$ of $E$ into a separable closure, the equation
\[
 \sum_j a_j\bigl(\tau(r_j)-\tau'(r_j)\bigr)=0
\]
defines a proper $k$-linear subspace of $k^N$. Finitely many such proper subspaces do not cover $k^N$. Choose the coefficients outside their union and set $t=\sum_j a_js_j$. The morphism
\[
 \nu=[s_0:\cdots:s_d:t]:X\longrightarrow\mathbf P_k^{d+1}
\]
is everywhere defined, has pullback $\mathcal O(1)$ equal to $L$, and is finite by the same fibre argument. Its scheme-theoretic image $Y$ is integral of dimension $d$. A homogeneous height-one prime in the polynomial coordinate ring of $\mathbf P^{d+1}$ is principal, so $Y$ is a hypersurface. The primitive-element choice gives $k(Y)=F(r)=E$, hence $\nu$ is birational.

For an affine open $\operatorname{Spec}A\subset Y$, write its inverse image as $\operatorname{Spec}B$. The ring $B$ is finite over $A$, has fraction field $E$, and is integrally closed because $X$ is smooth. Thus $B$ is the integral closure of $A$ in $E$: every element integral over $A$ is also integral over $B$, hence belongs to $B$, and the reverse inclusion follows from finiteness. This proves the normalization assertion with its exact map; compare the integral-closure argument of \cite[Tag 0AB1, Lemma 29.55.8]{AuditStacksFiniteNormalization}.

For the cohomological boundary, take
\[
 Y=\{y^2z=x^2(x+z)\}\subset\mathbf P^2_\C,
 \qquad
 \nu([u:v])=[v(u^2-v^2):u(u^2-v^2):v^3].
\]
The formula has no common zero and parametrizes $Y$ birationally. It is its normalization, with the two points $[1:1]$ and $[-1:1]$ over the node $P=[0:0:1]$. On the complex analytic spaces there is an exact sequence
\[
 0\longrightarrow\Q_Y\longrightarrow\nu_*\Q_{\mathbf P^1}
 \xrightarrow{\delta}\Q_P\longrightarrow0,
 \qquad \delta(a,b)=a-b
\]
at the node; away from the node the first map is an isomorphism. Since $\nu$ is finite, its higher direct images on these constant sheaves vanish. The map from global constants to $H^0(Y,\Q_P)=\Q$ is zero. The long exact sequence therefore gives
\[
 H^0(Y,\Q_P)=\Q\xrightarrow{\sim}H^1(Y(\C),\Q),
 \qquad
 \nu^*:H^1(Y(\C),\Q)\longrightarrow H^1(\mathbf P^1(\C),\Q)=0.
\]
Thus the geometric reduction comes with specified normalization data; it does not alone provide an ordinary-cohomology identification for the singular image.

The original cote14 scans111--113 (archive pp110--112) contain the global cycle-Lefschetz restrictions, Gysin maps, and punctured-cone formulations; scan131 (archive p130) contains a generic-member/base-locus diagram and its cycle-image formula. These inspected pages do not yet identify a proof carrying out the particular singular-hypersurface proposal in the 1967 Coates letter. I.M05.01 and II.MB08.01 retain that exact documentary obligation. The present construction is independently proved here, not attributed to an unseen manuscript passage.
\end{proof}
